
\subsubsection{3.1 Experimental Design}

The central methodological challenge is disentangling mechanism from computational budget. A comparison of semantic entropy (K=10 samples) to token variance (K=1 sample) confounds the clustering mechanism with sampling cost. We address this through matched ablation: both semantic entropy and the ensemble baseline use K=10 samples, isolating the clustering contribution.

For method independence analysis, we apply all methods to identical samples from the same model runs, ensuring that observed correlations reflect mechanistic differences rather than sample variance.

\subsubsection{3.2 Uncertainty Methods}

\textbf{Semantic Entropy} follows Kuhn et al. (2023). For each question, we sample K=10 answers at temperature 0.7. Answers are embedded using sentence-transformers (all-MiniLM-L6-v2) and clustered via agglomerative clustering with cosine distance threshold 0.5. This groups semantically equivalent answers. Entropy is computed over the cluster distribution: H_semantic = -∑_c p(c) log p(c), where p(c) is the proportion of samples in cluster c.

\textbf{Ensemble Baseline} provides the ablation control. We sample K=10 answers at temperature 0.7 (identical to semantic entropy) but use exact string matching instead of semantic clustering. Disagreement rate is U_ensemble = 1 - (count of most common answer / K). Comparing semantic entropy to this baseline isolates the clustering contribution.

\textbf{Self-Consistency} follows Wang et al. (2022). We sample K=10 answers and compute the agreement fraction for the most common answer.

\textbf{Verbalized Confidence} prompts the model to self-report confidence: "Provide your answer and your confidence level from 0\% (not confident) to 100\% (fully confident)." Confidence percentages are extracted via regex, with fallback to 50\% if no percentage is found. This requires K=2 forward passes.

\textbf{Token Variance} was intended to compute variance over token probability distributions. Due to an implementation issue discovered during validation, this method produced identical outputs to self-consistency (correlation 1.0) and requires reimplementation.

\subsubsection{3.3 Datasets}

\textbf{NaturalQuestions} provides factual questions from Google search queries. We used the unanswerable subset (100 questions) where no answer exists in the provided context. These questions were hypothesized to represent knowledge gaps.

\textbf{TruthfulQA} tests whether models generate truthful answers to questions designed to elicit common misconceptions. We sampled 100 questions. These questions were hypothesized to represent scenarios where models might confidently produce false information.

The hypothesis that these datasets cleanly separate error types was testable through diversity and agreement measurements.

\subsubsection{3.4 Model and Parameters}

All experiments used Mistral-7B-v0.1, an open-source 7-billion parameter decoder-only transformer. Generation parameters: temperature 0.7, maximum tokens 50, K=10 for semantic entropy/ensemble baseline/self-consistency, K=1 for token variance, K=2 for verbalized confidence, random seed 42.

\subsubsection{3.5 Evaluation}

\textbf{AUROC} (Area Under ROC Curve) measures discrimination between correct and incorrect answers, ranging from 0 to 1, with 0.5 indicating random performance.

\textbf{Pearson correlation} quantifies linear relationships between uncertainty method scores. Low correlation (|r| < 0.3) indicates orthogonal signals, high correlation (|r| > 0.7) suggests redundancy.

\textbf{Statistical significance} was assessed using t-tests for mean differences with significance threshold α = 0.05.
